\section{Method}
\label{sec:method}

Let \(Z_t\in\mathbb{R}^p\) be the causal state vector at sample \(t\). For compressor telemetry, \(Z_t\) contains the declared raw pressure, temperature, current, and flow channels used by the real-data diagnostic. For tabular condition-monitoring adapters, \(Z_t\) contains the declared feature vector for the cycle, vehicle, or production example.

Scaling is fitted on the training partition only. Each feature is centred by its training median and divided by a scale
\[
    s_j=\max\!\left\{\mathrm{IQR}_j,\ \kappa\left(q_{0.999,j}-q_{0.001,j}\right)\right\},\qquad \kappa=0.01,
\]
where all quantiles are training quantiles. The floor is not cosmetic. An interquartile range measures the spread of the central half, and a channel that idles at a fixed value between excursions has a central half made of quantisation noise: MetroPT's \texttt{pressure\_tp2} has an interquartile range of \(0.004\) against a \(0.1\)-to-\(99.9\) percentile range near ten, so its central half occupies four hundredths of one percent of its range. Dividing by that spread multiplies ordinary variation by roughly two hundred and fifty, and in a maximum-absolute-\(z\) aggregate one such channel then sets the score, its threshold and its exceedances for every other channel. Flooring the denominator at a fraction of the feature's own range leaves well-behaved features untouched, since for them the interquartile range already exceeds \(\kappa\) times the range, and prevents a quantised or mostly idle channel from setting the scale for the rest. The range is measured between the \(0.1\)st and \(99.9\)th percentiles rather than the 1st and 99th because a channel whose excursions are rarer than one percent has a narrow 1-to-99 range as well, which would leave the floor as degenerate as the quantity it bounds.

Operating regimes \(R_t\), when used in the sensitivity assets, are inferred without test labels. Regimes are conditioning variables for threshold estimation and are not interpreted as mechanical causes. The registered real-data matrix reported in Table~\ref{tab:event-benchmark-compact} and Supplementary Table~\ref{X-tab:non-event-results} uses a conservative train-split robust score rather than the full dangerous-region construction.

The implementation also provides dangerous-region scoring for separate diagnostics. The failure-proximal prototype variant uses only training failure intervals and records the source interval for every prototype. The unsupervised rare-state variant ignores failure labels and fits a high-distance rare-state set from training data only. These variants are available for diagnostic analysis and ablation coverage, but the registered real-data matrix in this manuscript is the conservative high-quantile score; it should not be read as evidence that both dangerous-region variants have already been optimized on every dataset.

For a target region \(A\), the scalar observable is
\[
    X_t=-\log\{d(Z_t,A)+\epsilon\},
\]
where \(d\) is the robust Euclidean distance in scaled coordinates and \(\epsilon>0\) prevents singularities. Larger \(X_t\) means closer proximity to the training-defined region. This is an anomaly or recurrence score, not a calibrated probability of failure.

For each regime \(r\), a threshold \(u_r(q)\) is estimated from the declared training or calibration partition at quantile \(q\). Exceedances are \(I_t(q,r)=\mathbf{1}\{X_t>u_r(q),R_t=r\}\). Runs declustering merges exceedances separated by at most \(m\) samples. The runs estimator, Ferro--Segers inter-exceedance estimator, and K-gaps-style declustering diagnostic summarize clustering under the frozen \(q\) and \(m\) settings \citep{FerroSegers2003Clusters,SuvegesDavison2010KGap}. The inverse extremal index is discussed only as an asymptotic cluster-size diagnostic under the assumptions of the estimator.

Pointwise exceedances are converted into alarm episodes by the declared merge gap. A test alarm is a true event alarm only if the episode intersects the predeclared warning window for a held-out failure under the frozen matching tolerance. All event recall, event precision, false-alarm burden, duplicate-alarm count, lead time, time under warning, and utility metrics are computed after this episode conversion when event labels exist. This prevents timestamp counts from being treated as independent maintenance evidence.

\paragraph{What is being evaluated.} Table~\ref{tab:method-taxonomy} states, for each of the nineteen benchmarked methods, what its score measures, how its threshold is chosen, how exceedances become episodes, which partition its parameters are fitted on, and where extreme-value theory enters. Three facts in that table matter for how this paper should be read, and none of them is apparent from the method names.

First, the recurrence observable \(X_t=-\log\{d(Z_t,A)+\epsilon\}\) defined above is used by the three target-region methods. It is not the score of the method labelled as registered. That method scores a robust maximum absolute \(z\), which is a multivariate anomaly magnitude, and it shares that score with seven of the baselines; what distinguishes it from them is its event policy.

Second, extreme-value theory changes the alarms a method raises in only five of the nineteen. A tail model sets the threshold for the peaks-over-threshold baseline and for SPOT, and an estimated extremal index sets the declustering run length for the Ferro--Segers policy, the K-gaps policy and the registered method. For the remaining fourteen the extremal-index estimates reported here are diagnostics computed beside the alarms; those methods would raise identical alarms if the estimates were discarded.

Third, and following from the first two, the contribution offered here is the evaluation architecture rather than a detector. What this paper supplies is a layered decomposition that locates where a detector loses its signal, a rank separation and a shift-based null that say whether a match carries information, and an event-level protocol under which simple and elaborate methods are compared on the same terms. The dynamical-EVT material motivates the recurrence observable and supplies the declustering rule, and the empirical results are evidence about what those produce under this protocol on two failures. They are not evidence that a robust anomaly score thresholded and declustered constitutes a dynamical-EVT detector, and the paper does not claim it.

\input{../reports/paper/latex/method_taxonomy.tex}

\paragraph{Warning-window rank separation and its null.} To separate a score that carries no pre-onset ordering from a threshold that discarded the ordering a score did carry, we compute the probability that a sample drawn from the declared horizon before a failure onset outranks one drawn from the rest of the test split. This is the Mann--Whitney statistic on midranks, so it is invariant to any monotone transformation of the score and comparable across methods; \(0.5\) is chance. Midranks matter here because industrial telemetry ties heavily during idle operation, and breaking ties by position would make a constant score appear informative purely because the warning window sits at one end of the series.

A recall of one is only evidence of anticipation if a detector could not have achieved it by alarming often enough. We therefore compute a chance-match probability by circularly shifting the entire alarm series by a uniformly random offset and asking how often at least one episode still intersects a matching window. The shift preserves the episode count, the episode durations, and the burstiness and regime structure of their spacing, and randomises only the alignment to the failure. Because an episode \([s,s+L]\) intersects a window \([a,b]\) exactly when \((s+d)\bmod n\) falls in an interval of length \(b-a+L+1\), the quantity is a measure of a union of circular intervals and is computed exactly rather than sampled.

We do not use a homogeneous Poisson model for this null. Poisson arrivals are independent and occur at a constant rate, whereas these episodes are produced by declustering a temporally dependent score, their rate varies with operating regime, and the merge rule induces further dependence between neighbouring episodes. A null that assumes extremes do not cluster cannot calibrate a paper about clustered extremes.

\paragraph{Event matching.} Let \(P=(p_1,\dots,p_m)\) be the alarm episodes after merging and \(F=(f_1,\dots,f_k)\) the labelled failures. For failure \(f\) with onset \(f^{\mathrm{start}}\), the matching window under horizon \(h\), pre-tolerance \(\tau_-\) and post-tolerance \(\tau_+\) is
\[
    W(f)=\left[\max\{0,\ f^{\mathrm{start}}-h-\tau_-\},\ f^{\mathrm{start}}+\tau_+\right].
\]
The window ends at onset plus \(\tau_+\), not at the end of the failure interval, so an alarm raised well inside a multi-day failure is not counted as a warning. An alarm \(p\) is \emph{eligible} for \(f\) when \(p\cap W(f)\neq\emptyset\); an alarm that begins before the window and extends into it is eligible, since only intersection is required. Matching is one-to-one: each alarm may be assigned to at most one failure and each failure to at most one alarm, chosen by a maximum-cardinality assignment over the eligible pairs, with ties broken toward the longer lead time. Alarm reuse across failures is available in the protocol but disabled here. Every alarm then falls in exactly one of three classes:
\[
\begin{array}{ll}
\text{matched} & \text{assigned to some failure;}\\
\text{duplicate} & \text{eligible for some failure but not assigned;}\\
\text{false} & \text{eligible for no failure.}
\end{array}
\]
Event recall is matched failures over \(|F|\). Event precision is
\[
    \mathrm{precision}=\frac{\#\text{matched}}{\#\text{matched}+\#\text{duplicate}+\#\text{false}}=\frac{\#\text{matched}}{m},
\]
so duplicates \emph{do} enter the denominator and precision is matched alarms over all alarms raised. This is deliberate: an operator responds to a duplicate exactly as to any other alarm, and excluding duplicates would let a detector improve its precision by emitting more alarms near a failure it has already matched.

These conventions determine the reported numbers, and one consequence deserves emphasis. Because episodes are formed by merging before matching, a detector that is permanently on produces a single episode that is matched, with no duplicates and no false alarms, and therefore attains precision and recall of one. Episode-level metrics cannot detect that state, which is why every alarm-burden comparison in this paper reports time under warning alongside episode counts, and why the benchmark flags any method whose warning exposure exceeds half the test period.

\paragraph{Joint utility.} The control comparisons use one scalar, declared before the comparison that consumes it. For a detector with detection indicator \(D\in\{0,1\}\), false alarms per operating day \(a\), warning exposure fraction \(c\), and lead time \(\ell\) in samples,
\[
    U=w_D D-w_a a-w_c c+w_\ell D\,\frac{\min\{\ell,\ell_{\max}\}}{\ell_{\max}},
\]
with \(w_D=100\), \(w_a=1\), \(w_c=10\), \(w_\ell=5\) and \(\ell_{\max}=7200\) samples. The units are therefore: one detection is worth 100, one false alarm per operating day costs 1, being under warning for the whole test period costs 10, and a lead time at or beyond two hours adds 5. The ordering is the substantive claim. Detection dominates by construction, so a method must avoid more than a hundred false alarms per day to justify missing a failure, which no method here approaches; the lead-time reward is capped and paid only on detection, so it cannot make a missed failure look competitive. Section~\ref{sec:results} reports a sensitivity sweep varying each weight by factors from \(0.25\) to \(4\), and the component metrics are reported separately so a reader who prefers different weights can recompute.

The score-stratification target for horizon \(h\) is
\[
    Y_t^{(h)}=\mathbf{1}\{\hbox{a labelled event begins within }h\hbox{ after }t\}.
\]
In the current evidence, model outputs are anomaly ranks or recurrence scores. Probability-calibration language is therefore restricted: Brier-style summaries are reported only as diagnostics of score stratification unless a separate calibration partition contains enough positive independent units for an explicit probability map.

The leakage contract is: training features cannot contain held-out failure samples; scaling, regimes, target regions, thresholds, lag selection, merge-gap selection, and calibration cannot use test events; SCANIA vehicles cannot appear in multiple partitions; and feature windows cannot cross maintenance-reset boundaries. Table~\ref{X-tab:leakage-audit} gives the terminal audit result for each registered check. Under the current registered diagnostic the algorithm is:
\[
\begin{array}{ll}
1. & \text{Split entities or intervals into training, validation, calibration, and test partitions.}\\
2. & \text{Fit scaling and robust score thresholds on allowed partitions only.}\\
3. & \text{Select } q,m,\text{ and merge gap on validation/calibration partitions only where used.}\\
4. & \text{Freeze all parameters and score the held-out test partition causally.}\\
5. & \text{Extract exceedance clusters, merge alarm episodes, and match them to labelled events.}\\
6. & \text{Report event metrics, uncertainty by independent unit, and non-estimable cases.}
\end{array}
\]


The practical complexity for scoring is \(O(npK)\), where \(n\) is the number of samples or entities, \(p\) is the state dimension, and \(K\) is the number of stored prototypes or rare-state summaries. Thresholding and runs declustering are linear in \(n\). The streaming-compatible parts are causal feature construction, scaling with frozen parameters, scoring, thresholding, and alarm merging. Prototype construction, model selection, and final event matching are offline evaluation steps.
